\documentclass[11pt,a4paper]{article}
\usepackage[margin=0.9in]{geometry}
\usepackage{amsmath,amssymb,booktabs}
\setlength{\parindent}{0pt}\setlength{\parskip}{0.7em}
\begin{document}

\begin{center}
{\Large\textbf{Mentor review package}}\\[0.4em]
\textbf{HUMAN ACTION REQUIRED} --- ten items, G-1 through G-10.\\
Every one is a \emph{confirmation} of something already computed. None asks for a calculation.
\end{center}

Decision form: \texttt{review/MENTOR\_DECISION\_FORM.md}.\\
Full item text: \texttt{spinor\_trace\_bridge/docs/MENTOR\_REVIEW\_ITEMS.md}.

Read in this order. The first three are the ones where a wrong answer changes a
number rather than a sentence.

\hrulefill

\section*{Priority 1 --- G-10: the leading-degree rule}

\textbf{This is the only unverified analytic input the degree-ten result rests on.}

The stress-flow closure assigns each generated target to a graded piece using a
leading field degree per generator. The rule:
\begin{align*}
\mathrm{Tr}(\tau)   &\;\longrightarrow\; 4, \text{ not } 2\\
\mathrm{Tr}(\tau^k) &\;\longrightarrow\; 2k, \quad k \ge 2\\
\text{products}     &\;\longrightarrow\; \text{additive}
\end{align*}
All 18 generators at all six primes place their first appearance exactly where
this predicts (\texttt{tests/test\_leading\_degree\_rule.py}), so the assignments
are no longer self-referential --- they agree with a rule written down separately.

\textbf{G-10 has since been derived}, not merely checked against the table. For
any improvement coefficient $c$ the trace is $(1-cd)\langle F,F\rangle$, so only
$\langle F,F\rangle$ matters; $F\wedge F$ is a top form built from two copies of
an odd-degree form and vanishes; $\langle F,\star F\rangle$ is proportional to
it; hence $\langle F,F\rangle=0$ on either eigenspace. Verified with a control at
four primes by code importing no flow machinery, and load-bearing: forcing a
degree-2 contribution gives $\dim Q_{10}=0$ instead of $3$.

\textbf{What only you can confirm} is now narrower: that the free theory and
stress convention are the ones you intend, and that no formulation in use changes
the \emph{quadratic scalar available}. Independence from $c$ means a different
trace convention alone cannot change the conclusion.

\textbf{If this is wrong:} every target lands in the wrong graded piece, and
$\dim D_{10} = 11$ --- hence $\dim Q_{10} = 3$, the paper's central number --- is
wrong with it. Nothing else in the paper has this property.

\section*{Priority 2 --- G-2: the split-signature correction}

An earlier record in this project stated the oscillator frame's real form is
Euclidean $SO(10)$, where $\star^2 = -1$ admits no real self-dual five-form, and
treated that as blocking the bridge.

\textbf{That was wrong.} Computing the metric from the archive's own wedge and
contraction operators --- all one hundred anticommutators, diagonalised --- gives
eigenvalues $(+\tfrac12)^5, (-\tfrac12)^5$: real signature $(5,5)$, split. A null
frame cannot be Euclidean at all, since Euclidean signature has no isotropic
vectors.

\begin{center}
\begin{tabular}{llll}
\toprule
 & real form & $\star^2$ on 5-forms & real self-dual 5-forms \\
\midrule
Lorentzian   & $\mathrm{Spin}(1,9)$            & $+1$ & yes \\
split        & $\mathrm{Spin}(5,5)$            & $+1$ & yes \\
Euclidean    & $\mathrm{Spin}(10)$             & $-1$ & \textbf{no} \\
complexified & $\mathrm{Spin}(10,\mathbb{C})$  & ---  & the common ground \\
\bottomrule
\end{tabular}
\end{center}

What survives is weaker and precise: $(5,5)$ and $(1,9)$ are inequivalent
\textbf{real} forms, so a real frame transformation between them does not exist.
But both metrics have discriminant $-1$ up to squares, so over $\mathbb{C}$ and
over $\mathbb{F}_p$ they are congruent, and the bridge constructs that congruence
explicitly and checks it. Every component-level comparison in the paper is
therefore modular, where the transition exists exactly --- never a real Lorentzian
identification.

\textbf{What is needed:} confirmation that the correction is accepted, and ---
because the earlier, wrong statement may already have been communicated to you ---
that it has not been relied on elsewhere.

\section*{Priority 3 --- G-8: credit for rank 81}

The number 81 is \textbf{not ours} and the paper says so.

\begin{center}
\begin{tabular}{ll}
\toprule
element & origin \\
\midrule
analytic count $126-45=81$ from the trivial generic stabiliser & \textbf{literature}, cited \\
earlier float64 finite-difference Jacobian evidence & \textbf{mentor archive} \\
exact analytic amputated derivative, no step size, no tolerance & this project \\
integral basis making the modular rank a rigorous bound & this project \\
explicit $81\times81$ minor, two determinant routines & this project \\
complete 83/83 candidate schedule with terminal statuses & this project \\
\bottomrule
\end{tabular}
\end{center}

The claim split in the draft: the \emph{upper} bound is analytic and attributed;
only the \emph{lower} bound $\mathrm{rank}_{\mathbb{Q}} \ge 81$ is claimed here,
and only at finitely many sample points. A wording gate fails the build on
``proved rank 81 computationally''.

\textbf{What is needed:} confirm the division of credit and that the literature
attribution points at the source you intend. Also confirm you want the paper to
state plainly that 83 candidates of rank 81 means the selection carries
functional dependencies --- it currently does.

\section*{Priority 4 --- G-9: how to present the cardinality proposition}

Any set closing degree $d$ has at least $\dim Q_d$ elements, because the quotient
map is linear and a span of $k$ vectors has dimension at most $k$. Three lines.

It is in the paper because it discharges half of PO-08 and turns a
basis-dependent assertion into a basis-free one --- not because it is new. It is
almost certainly not new as mathematics, and the draft makes no novelty claim.

\textbf{Options:} proposition with proof (current) $\cdot$ lemma $\cdot$ remark
$\cdot$ corollary $\cdot$ omit as elementary.

\section*{Remaining items}

\begin{center}
\begin{tabular}{lll}
\toprule
item & subject & what a wrong answer costs \\
\midrule
G-1 & reconstructed spinor index placement & a wording change \\
G-3 & real-form independence of dimensions & a citation instead of an argument \\
G-4 & equivariance at sampled group elements & possibly a symbolic proof \\
G-5 & comparison is modular, not char.\ zero & wording only \\
G-6 & real-form caveat; AMB-01/02 source intent & wording; no result depends on it \\
G-7 & redistribution of the spinor archive & release contents only \\
\bottomrule
\end{tabular}
\end{center}

\section*{What is not being asked}

No item asks you to approve a submission, an author list, or a novelty claim.
Those are separate and are in \texttt{audit/AUTHORSHIP\_AND\_CREDIT\_FINAL.md}.
Every novelty row is \texttt{PROVISIONAL} and stays that way until someone who
knows the field says otherwise --- a literature search shows absence of evidence,
not evidence of absence.

\end{document}
